\section{Related Work}

\subsection{Co-Training for Robot Learning}
Generally, we can define any learning system that utilizes heterogeneous data as co-training.
To mitigate the gap between limited in-domain robot data and large-scale surrogate data or even multimodal resources, co-training has been employed in numerous studies. 
Based on the types of large-scale surrogate data, we can coarsely categorize current work into the following three intersecting types: \textit{sim-and-real co-training}, \textit{cross-embodiment co-training}, and \textit{non-robot data co-training}.

\textbf{Sim-and-Real Co-Training.}
Using simulation as a surrogate data source is a promising way for co-training.
The main domain gaps include visual and physics gaps, with the physics gap being much more challenging.
But with the advancement of high-fidelity physics simulators ~\cite{mittal2025isaac, nasiriany2024robocasa, zakka2025mujoco, Genesis} and automated data generation tools ~\cite{mandlekar2023mimicgen, jiang2025dexmimicgen, lin2025constraint}, high-quality and massive robot trajectories can be obtained easily.
Many works have demonstrated the effectiveness of sim-and-real co-training on challenging manipulation tasks with small diffusion policy models ~\cite{maddukuri2025sim, wei2025empirical, cheng2025generalizable} and even large Vision-Language-Action(VLA) models ~\cite{bjorck2025gr00t, yu2025real2render2real}.
There are also works ~\cite{barreiros2025careful, jain2025polaris} utilizing relative large-scale real-world data and small amount of simulation data for co-training, so that the performance evaluated in simulation can effectively reflect the performance in real world.


\textbf{Cross-Embodiment Co-Training.}
A large amount of works ~\cite{doshi2024scaling, yang2024pushing, o2024open, intelligence2025pi_, yuan2025motiontrans} have explored using cross-embodiment robot data for co-training, where a single policy is trained with multiple embodiments with a unified architecture and action representation.
The main domain gaps include the embodiment appearance and the embodiment-physics gap.
Human data is a special case of them, which can be directly treated as another embodiment.
These works show that diverse robot pretraining can produce transferable internal representations ~\cite{kareer2025emergence}.
Some works ~\cite{cai2025n, kareer2025egomimic, punamiya2025egobridge} utilize representation alignment regularization like optimal transport or adversarial discriminator.
Another line of works ~\cite{xu2025dexumi, lepert2025phantom, lepert2025masquerade} force input data distribution alignment via image editing.


\textbf{Non-Robot Data Co-Training.}
Internet-scale multimodal data without action label is another valuable co-training resource.
Numerous works ~\cite{intelligence2025pi_, lee2025molmoact, zawalski2024robotic, zitkovich2023rt, cotraininglbm2025} have adopted VL datasets, which contain rich commonsense knowledge, planning and spatial information for co-training in VLAs.
These works have demonstrated that co-training with VLM data can transfer knowledge from other modalities.
Also a number of works try to utilize videos for policy co-training.
Some works ~\cite{kareer2025egomimic, lepert2025masquerade, luo2025being} explicitly extract action labels but with the problem of accuracy;
some works ~\cite{ye2024latent, zhou2024dino} explores latent action representations by encoding changes between video frames;
other works ~\cite{li2025unified, zhu2025unified, liang2025video} propose to co-train the action- and actionless video data within a unified architecture.


\subsection{Representation Learning in Domain Adaptation}
Co-training can be also viewed as semi-supervised or unsupervised domain adaptation.
A central theme in domain adaptation is learning representations that enable knowledge transfer across domains while mitigating distribution shift. 
Early theoretical work formalized this goal by relating target error to source error and representation-level domain discrepancy, motivating the pursuit of domain-invariant features through shared embeddings or feature transformations~\cite{ben2010theory, mansour2009domain}.    
Building on this foundation, many practical approaches explicitly align representations across domains, including discrepancy-based methods that minimize statistical distances such as maximum mean discrepancy (MMD)~\cite{long2015learning}, OT-based alignment~\cite{courty2017joint, damodaran2018deepjdot}, and adversarial domain adaptation methods that encourage indistinguishability ~\cite{ganin2016domain, tzeng2017adversarial}. 
